\paragraph{Attaching a block of sources.}
Fix $K\geq3$. Let $B$ be core on a nonempty finite set $J$ of inputs, let $W$ be a
disjoint nonempty finite set, and write $U_W(w)=\bigcup_{i\in W}w_i$.
For a proper coloring $q$ on $J$, let $L_q$ consist of the background
valuations satisfying every full factor of $\Theta_{B,q}$. These
valuations may use every singleton color. Given any upper bound $F$
with $B\subseteq F\subseteq B+U_W$, let $D_q$ be the union of the
core operations $h$ on $W$ satisfying
\[
 h(w)\subseteq F(u,w)\qquad(u\in L_q,\ \text{all } w).
\]
It is the greatest such operation, or $0$ for an empty union. Give
$0$ the constantly true principal. Write $E(u)^{\widehat k}$ for the
singleton coloring of $u$ with the $k$-block erased.

\begin{theorem}[Block extension and activation]
\label{core:thm:block-source-extension}
The core operations between $B$ and $B+U_W$ are exactly
\[
 B(u)\ \cup\ \{k:k\in h_{E(u)^{\widehat k}}(w)\},
 \qquad p\preceq_B q\ \Longrightarrow\ h_q\subseteq h_p ,
\]
where each $h_q$ is a core operation on $W$ or zero. For a core
extension $F$, the $D_q$ form its greatest representing family.
For an arbitrary upper bound $F$ as above, the greatest core
operation below it is
\[
 B(u)\ \cup\ \{k:k\in D_{E(u)^{\widehat k}}(w)\}.
\]
In particular $F$ is core precisely when it equals this operation.
When $F$ is core, its entire activation preorder is given by
\[
 p\preceq_F q
 \quad\Longleftrightarrow\quad
 p_J\preceq_B q_J\quad\text{and}\quad
 \Theta_{D_{q_J},q_W}(p_W)=1
\]
for all proper colorings on $J\cup W$, with subscripts denoting
restriction. No monotonicity of the proposed upper bound is required.
\end{theorem}
\begin{proof}
For a word $t$ write $\lambda_i^t(d)$ for survival of source $i$
with the other inputs colored by $d$, even if $d$ uses every owner.
To read this bit, give $i$ a fresh singleton $\{\omega\}$, use the
specified singleton menus, give the holes the ambiguous menu $I$,
and test for $\omega$. Since no gate is owned by $\omega$, a gate
can lose the mark only to its owner's singleton in the other branch.
The test thus reads the full path factors of the occurrences of $i$
in the chosen word. When $d$ omits an
ordinary owner $k$, the bit is the ordinary $k$-source coefficient.
Indeed the foreign singleton tests are the same, and a $k$-gate
passes membership of $k$ by disjunction, as it does for the mark
when no input is the $k$ singleton.

Fix a background coloring $e$. From $t$ obtain a word $h_e$ on $W$,
possibly zero, by fixing and erasing the colored inputs one at a time
using \cref{core:lem:winner-erasure}, and neutrally erasing the holes.
Use the carrier with the extra winner for these erasures. The
compilers use original owned gates and union, which can be expanded
at three original owners. Each erasure preserves $\omega$ at the
specified fixing, since the fixed menu omits it. A constant fixed
singleton or an empty residue gives zero for the remaining
bits. Each nonzero erasure is bounded by its original word at every
ambiguous menu in that coordinate, by the neutral-residue bound.
Thus, simultaneously at every $i\in W$ and coloring $d$ on
$W\setminus\{i\}$,
\[
 \lambda_i^{h_e}(d)=\lambda_i^t(e,d).
\]
Zero has no surviving sources. Only word evaluations use the extra
winner; all bounds below concern the original menus.

Suppose $t\subseteq B+U_W$. For fixed $i,d$, the section
$e\mapsto\lambda_i^t(e,d)$ on proper $J$-colorings is an actual
interval predicate above $B$. Indeed erase $W\setminus\{i\}$ at $d$
by the same construction. Its word $g$ is below $B+v_i$: to test a
winner $k$, give every erased input the ambiguous menu
$I\setminus\{k\}$. Then $B+g$ is an interval realizer with the
displayed predicate, also when $g=0$.

It follows that $h_e(w)\subseteq t(u,w)$ for every $u\in L_e$.
For $k\in h_e(w)$ choose a supplying $i$ and take $d$ to be the other
$W$-statuses with their $k$-block erased. The section is positive at
$e$. It is therefore positive at $E(u)^{\widehat k}\preceq_B e$.
This last coloring joined with $d$ omits $k$, so the ordinary source
test gives $k\in t(u,w)$.

Now take any core $t\subseteq F$. If $k\in t(u,w)\setminus B(u)$,
some $W$-source supplies it. A $J$-source bit of $t$ is bounded by
that of $B$, as its sole-source probe makes $U_W$ omit the queried
winner. For $e=E(u)^{\widehat k}$ the supplying $W$-bit survives in
$h_e$, which lies below $D_e$ by the preceding containment. This
proves the claimed upper bound for every $t$.
To attain it, for each proper $e$ take a one-source principal spine
$S_e$ above $B$ and substitute $D_e(w)$ for its source, omitting
zero bodies. Take their union with $B$. A nonconstant cofactor of
$S_e$ is union with $B(u)$ and requires $u\in L_e$, since an
absorbing path helper would make it constant. Thus every summand
is bounded by $F$. For $e=E(u)^{\widehat k}$ its predicate is
positive at $e$, so the summand supplies $k$ whenever $D_e(w)$ does.
This proves the greatest-minorant formula.

The sets $L_e$ increase in the activation preorder, by testing each
monochromatic block, so $D_e$ decreases. Conversely use any antitone
family $h_e$ in the same construction. At $p=E(u)^{\widehat k}$ a
summand can supply $k$ from its body only if $p\preceq_B e$, whence
$h_e\subseteq h_p$, and the summand at $p$ supplies every such
$k\in h_p(w)$. Thus the resulting operation has the asserted form.
On $L_e$ every winner-erased coloring lies below $e$, so $h_e$ is
bounded there by the operation it represents. This gives $h_e\subseteq D_e$.

Finally let $F$ be core. Every $F$-comparison restricts to a
$B$-comparison by adjoining $F$ to any $B$-interval realizer.
Put $e=q_J$. For $D_e\ne0$ and $G\in\mathcal I(D_e)$, the term
$F+S_e(G(Y,w),u)$ belongs to $\mathcal I(F)$ by the same cofactor
bound. Its source predicate is the product of $\Theta_{B,e}$ and
the predicate of $G$. This proves necessity of the second comparison.
Conversely take a principal spine $S$ above $F$ positive at $q$.
Apply the erasure at $e$ to $S$, with $Y$ included in the block.
The resulting $h_e(Y,w)$ has the exact formal section of its
$Y$-source bits and is bounded by $F+Y$ on $L_e$, since
$S\subseteq B+U_W+Y$. Hence $\partial_Y h_e\subseteq D_e$, by
testing each winner with $Y=I\setminus\{k\}$. Ghost comparison gives
$h_e\subseteq D_e+Y$. Its section is positive at $q_W$, so $h_e$
is nonzero. If $D_e=0$ the word is just $Y$, since any other
occurrence gives nonempty residue. Otherwise $D_e+h_e$ is an interval
realizer above $D_e$. In both cases the second comparison gives
positivity at $(q_J,p_W)$ in the formal section of $S$.
Fixing $p_W$ gives an actual $B$-interval predicate in the first
block, by the section argument with $Y$ marked. The first comparison
therefore transfers positivity to $p$, as required.
\end{proof}

\paragraph{Reading the coefficients from guards.}
Suppose the sources in $J$ have the predicates of $B$ independently
of $W$. For $i\in W$, write $(b,t)$ for a minimal proper coloring
at which its proposed antitone source predicate $\pi_i$ vanishes.
Here $b$ colors inputs in $J$ and $t$ colors inputs in
$W\setminus\{i\}$. For each proper $q$ on $J$ set
\[
 \mathcal T_{q,i}=\min_{\subseteq}
       \{t:(b,t)\text{ is a minimal zero of }\pi_i,\ b\preceq_B q\}.
\]

\begin{lemma}[Restriction of minimal blocking patterns]
\label{core:lem:block-source-patterns}
For the candidate reconstructed from these predicates, $D_q$ is the
greatest core operation $h$ on $W$ satisfying
$\pi_i^h(t)=0$ for every $i\in W$ and $t\in\mathcal T_{q,i}$,
or zero if there is none. Give zero the constantly false source
predicates. The candidate is core exactly when
\[
     \pi_i(e,d)=\pi_i^{D_e}(d)
     \qquad(i\in W,\ (e,d)\text{ proper}).
\]
\end{lemma}
\begin{proof}
For necessity choose an owner $k$ missing from $(b,t)$, use the
specified atoms, set $i$ to $\{k\}$ and every hole to
$I\setminus\{k\}$. The background lies in $L_q$ since its full
factors at $b$ hold. Only $i$ could supply $k$, and the proposed
guard blocks it. Thus every bounded $h$ has a zero at $t$.
Conversely suppose $u\in L_q$ and a source $i$ supplies $k$ in
$h(w)$ outside the candidate. Put $e=E(u)^{\widehat k}$ and let
$d$ be the other internal statuses with their $k$-block erased.
Some minimal zero $(b,t)$ lies in $(e,d)$. Since
$b\subseteq e\preceq_B q$, the constraints and antitonicity give
$\pi_i^h(d)=0$, a contradiction. Finally the greatest-minorant
formula in \cref{core:thm:block-source-extension} has exactly the
displayed internal predicates, by the same sole-source probes, and
the prescribed background predicates. Source reconstruction proves
the criterion.
\end{proof}

\paragraph{How many sources can be lost?}
Call a source conditional if its predicate is not identically one;
this includes an unused source. When only these sources can be
lost, eliminating the retained inputs gives a direct implication
test on the minimal zeros. In particular the choice between the
two canonical presentations can matter already at two conditional
sources.

\begin{lemma}[Eliminating inputs outside the queried heads]
\label{core:lem:unmarked-background}
Consider one positive query $(i,e,d)$ and arbitrary negative queries
$(j,b,t)$. Let $H$ contain $i$ and all the negative heads, and
split the postings into their parts on $J=V\setminus H$ and $H$.
There is a witness exactly when a word on $H$ supplies $(i,d)$
and omits $(j,t)$ for each row with $b\subseteq e$.
\end{lemma}
\begin{proof}
The case $J=\varnothing$ is immediate. Form a complete upper bound
by retaining $J$ and prescribing just the upsets of the given
negatives as losses. It is nonempty because it contains the union
on $J$. A sample witness is below this bound by source decomposition
and antitonicity. Adjoining the union on $J$ preserves its
negative and positive bits. Conversely any word under the bound
supplying the positive is a witness. By
\cref{core:lem:block-source-patterns} the greatest coefficient at $e$
omits exactly the bits forced by the indicated internal constraints:
the activation preorder of the retained union is subcoloring by
\cref{core:thm:source-completion}. Thus its bit at $(i,d)$ is
positive precisely when some internal word supplies it. This
argument also covers an internal negative $t$ whose union with the
whole of $e$ uses every color: it suffices that its original row
$(b,t)$ is proper.
\end{proof}

Erasing background inputs also preserves a bound on the number of
completion layers. The singleton tests that cease to be attainable
under the fixing have to be recorded separately.

\begin{theorem}[Erasure at fixed completion depth]\label{core:thm:completion-erasure}
Fix the singleton statuses of the background ports $J$ of a completion
expression $t$, leaving ports $H$ free. A background menu $b_j$ with
neutral status may be any nonsingleton. Write $U_b=\bigcup_{j\in J}b_j$.
Either, for every such $b$, the output is constant in the free inputs
and lies in $U_b$, or there are one expression $s$ on $H$ of depth at
most that of $t$ and a set $M$ of owners such that for all $b,w$
\[
 s(w)\subseteq t(b,w)\subseteq s(w)\cup U_b,\qquad
 t(b,w)=\{r\}\ \Longleftrightarrow\ r\in M\text{ and }s(w)=\{r\}.
\]
In the constant case the singleton status depends only on the fixed
statuses and no free source survives. In the second case $s$ also preserves each free source's
survival at the fixing, including on colorings using all owners.
\end{theorem}
\begin{proof}
Proceed up the expression. A background port is constant; a free port
has mask $I$. At a free union keep the union of its surviving children
and intersect the sets of owners that can be singleton in its children,
using the singleton status of a constant child. If both children are
constant the union is constant.

At a gate let $P$ be the set of surviving child ports, with masks $M_j$.
The other children are constant. Let $C_r$ be the constant ports with
value $\{r\}$. On $A\subseteq P$ put
\[
 f_r(A)=\Phi_r\bigl(C_r\cup\{j\in A:r\in M_j\}\bigr).
\]
These functions give the original gate's atomic tests when $A$ lists
the surviving children whose residual values are $\{r\}$.
If $f_r(\varnothing)=1$, the output is always $\{r\}$, and some
constant port supplies $r$. With $P=\varnothing$ the output is also
constant. Otherwise put $M=\{r:f_r(P)=1\}$. On $P$ use $f_r$ for
$r\in M$ and unanimity for the other owners. For at least two active
owners a rainbow selection has a common port by the original profiles;
it cannot be constant since their $C_r$ are disjoint. For one active
owner its winning selection is nonempty since $f_r(\varnothing)=0$.
If none is selected, use any port in $P$. This gives a greatest
completion on the surviving children, hence preserves the depth bound.

An original singleton survives with its owner, and a new singleton
at an inactive owner $r$ requires every surviving residual to equal
$\{r\}$. In this last case the original gate takes the union of its
children, which contains $r$ and has no other winner outside $U_b$.
When both gates are ambiguous the two inclusions follow by taking
unions. This proves the assertions on ordinary menus. For source
survival, track a mark carried only by a free ambiguous input and
erased by any singleton output, as in the proof of
\cref{core:thm:block-source-extension}. At an inactive new singleton
none of the surviving children carries the mark. In every other case
the two gates both erase it or both pass it from their children.
This argument uses only the statuses, so no unused owner is needed.
\end{proof}

For a feasible sample in \cref{core:lem:unmarked-background}, let $d$
be its minimum completion depth and $d_H$ the minimum for the indicated
internal sample. Then
\[
                  d_H\ \le\ d\ \le\ d_H+\min(K,|J|).
\]
For the lower bound erase the background at $e$ in a witness of depth
$d$. Its surviving positive rules out a constant residue. On a negative
with $b\subseteq e$, extending $b$ to $e$ can only remove surviving
sources: on the path from a marked leaf, an absorbing sibling remains
absorbing when additional inputs are posted. Equivalently, at a
completion gate each positive atomic test persists under the extra
postings. The residual consequently meets all internal negatives.
For the upper bound start with an optimal internal expression.
Every other negative has some background posting $v:r$ with
$v\notin e_r$. Group these rows by such a color or by such an input.
As in \cref{core:thm:one-source-depth}, each group is enforced by one
greatest completion on the current witness and background projections.
The latter carry no marks at any negative head. This establishes the
inequality and, by using the minimizing partition of the background
tails in that theorem, can further reduce the overhead. The bound is
sharp already for $|H|=1$, where the internal depth is zero and a
matching of input--color pairs forces $\min(K,|J|)$ layers.

Two singleton fixings can force erasure to remove an entire layer.
Choose distinct $b,c$. For any nonempty-winner-set valued operation
$f$ introduce a fresh input $y$ and put
\[
 \widehat f(w,y)=
 \begin{cases}
  y,&y=\{b\}\text{ or }y=\{c\},\\
  f(w)\cup y,&\text{otherwise}.
 \end{cases}
\]

\begin{theorem}[Adding exactly one completion layer]\label{core:thm:completion-suspension}
For every $d\ge0$, $f$ has a realization of depth at most $d$ if and
only if $\widehat f$ has one of depth at most $d+1$.
No $\widehat f$ has depth zero. Thus this construction increases every
finite minimum depth by exactly one and preserves nonrealizability.
\end{theorem}
\begin{proof}
Apply a binary greatest completion to $f$ and $y$, with the second
projection's profiles at $b,c$ and unanimity elsewhere. This gives
the displayed operation.

Conversely, take a realization of $\widehat f$ and pass through outer
free unions to its summands. When $y=\{b\}$ or $y=\{c\}$ each summand
must take that same value, so its profiles at both owners win on the
single port $y$. By \cref{core:lem:singleton-profile-absorption} it
contains the whole menu of $y$ at every valuation: the two possible
absorbing singleton exceptions are distinct. On a neutral fixing of
$y$ every summand is therefore ambiguous. A projection summand can
only be $y$. Every outer completion summand now takes the union of its
children. Erase neutral $y$ from those children by
\cref{core:thm:completion-erasure} and collect their surviving
residuals in one free union. They each have depth at most $d$ and
depend only on the neutral status, independently of the menu chosen.
For any winner $k$, take $y=I\setminus\{k\}$, neutral since $K\ge3$.
The constant residues and the projection $y$ omit $k$, and a surviving
child contributes $k$ exactly when its residual does. The union thus
has exactly the memberships of $f$. There must be a surviving child
since $f$ has nonempty values. The same argument excludes depth zero,
where only projection summands would be available.
\end{proof}

On canonical source guards this transformation has linear cost at
fixed $K$, counting the listed guards and their literals. Give $y$ the
empty guard. At each old head and observer $k$, adjoin to every prime
the literals excluding the postings of $y$ in $\{b,c\}\setminus\{k\}$.
Every sufficient guard must exclude both visible postings, as can be
tested on $y$ alone. After these exclusions its other literals form
a sufficient guard exactly when they did before. This proves
canonicity, including empty lists at unused heads.

\begin{theorem}[Compatibility over retained inputs]
\label{core:thm:retained-background-cuts}
Suppose the family has a positive bit. Let $H$ be its nonempty set
of conditional sources and let $J$ be the retained inputs. Write a listed minimal zero as $(j,b,t)$,
where $b$ posts background inputs in $J$ and $t$ posts inputs in
$H\setminus\{j\}$. Fix $i\in H$, a proper internal coloring $d$ off
$i$, and an observer $k$ missing from $d$. Take the list of core
operations $u$ on $H$ for which $\pi_i^u(d)=1$. For each such $u$
select a listed zero $(j_u,b_u,t_u)$ with $\pi_{j_u}^u(t_u)=1$.
A complete family is realizable if and only if it obeys the following
rule for every such selection:
\[
 \text{if } b=\bigcup_u b_u\text{ is a consistent coloring omitting }k,
 \qquad \pi_i(b,d)=0.
\]
For fixed $K$ and $|H|$, this gives polynomial-time recognition
from the minimal zeros and coNP membership from the prime guards.
\end{theorem}
\begin{proof}
Use \cref{core:lem:unmarked-background} at a positive $(i,e,d)$
against all minimal zeros. Retained inputs supply their own
positives by projection. If all operations capable of supplying
$(i,d)$ are excluded, take one excluding row per operation.
Their backgrounds unite inside $e$. Conversely a consistent union
omitting $k$ gives a proper fixing at which all the selected operations
are excluded. Antitonicity proves the rule in both directions:
every positive must be supplied by the greatest allowed coefficient.
For a fixed internal arity the list of operations is constant.
With explicit minimal zeros one may enumerate the selections; with
prime guards one may guess proper zero rows instead of minimal ones,
since each contains a minimal zero, and test their bits and the
positive at their union directly.
\end{proof}

\begin{corollary}[A pair rule for two conditional sources]
\label{core:cor:two-conditional-guard-complexity}
Suppose $x,y$ are the only conditional sources and the family has
a positive bit. Choose a minimal zero at each of $x,y$. Remove
postings of $x,y$ from these two zeros and let $b$ be the union
of the remaining postings. If $b$ is a consistent proper coloring
and the two removed postings are not both present with the same
color, require
\[
                  \pi_x(b,\varnothing)=\pi_y(b,\varnothing)=0 .
\]
This one rule characterizes realization. Recognition from complete
prime guards is coNP-complete at each fixed $K\ge3$, even when
only one of $x,y$ depends on the other and every prime has at most
four literals. With just one conditional source and a retained
input every antitone observer-agreeing family is realized.
\end{corollary}
\begin{proof}
At a fixed proper background $e$, collect the internal constraints
whose backgrounds are contained in $e$, as in
\cref{core:thm:retained-background-cuts}. A constraint at $x$ rules
out the projection $x$, and likewise at $y$. If both heads have
constraints, a nonzero coefficient must therefore use both inputs.
By \cref{core:prop:binary} the only binary word losing each source
at some pattern is $x\star_r y$: in both directions the loss must
post the other input at the same $r$. An incompatible pair forces
coefficient zero already on its background union. This proves
necessity.

Conversely a background containing an incompatible pair has both
heads lost there already, by the rule and antitonicity.
If both heads have constraints but no such pair exists, all of them
post the same $r$ and $x\star_r y$ gives exactly the required bits
on the proper extensions of $e$. If only $x$ has constraints, use
$y$ when an empty internal pattern occurs. Otherwise retain $y$
and suppress $x$ just when $y$ has a color in the listed set, using
the corresponding one-sided binary normal form. Interchange the
heads for the other case and take union if neither has constraints.
The block criterion now realizes the family. The one-source
assertion is \cref{core:thm:source-completion}.

For the complexity lower bound, reduce failure of hypergraph
two-colorability, where each edge must meet both colors. 3SAT gives
such instances with edges of size at most four: use two complementary
vertices per variable joined by an edge, and a marker $z$. Each
clause gives the edge consisting of $z$ and its literal vertices.
Color $z$ false. The complementary pairs give a truth assignment
and the clause edges require a true literal. Remove edges containing
a smaller one and take a nonempty edge family without empty edges.

Choose distinct owners $a,b,c$ and retain one input per vertex.
Prescribe a zero at $x$ exactly when the vertices posted at $a$ meet
every edge. Prescribe a zero at $y$ exactly when $x$ is posted at
$c$ and the vertices posted at $b$ meet every edge. The primes at
$x$ for observers other than $a$ are the guards excluding the
$a$-postings of one whole edge. At $y$ for observers outside
$\{b,c\}$ they are the single exclusion of $x$ at $c$ and the
analogous edge guards for $b$. All remaining facets have the
empty guard. These lists are canonical: all allowed vertices of the
tested color can be posted together, so a sufficient guard must
disable the required internal posting or exclude every vertex of
one edge. The edge antichain gives minimality.

Every opposite-head pair has just one internal posting, at $x$.
Its backgrounds are compatible precisely when the chosen
$a$-transversal and $b$-transversal are disjoint. Their union omits
$c$ and leaves $y$ live without the internal posting. Hence the
pair rule fails exactly when two disjoint transversals exist.
These extend to a two-coloring; conversely each color class of a
two-coloring contains a minimal transversal. This proves the
lower bound. Membership follows from
\cref{core:thm:retained-background-cuts}.
\end{proof}

The pair rule is necessary even among arbitrarily many conditional
sources. Apply \cref{core:lem:unmarked-background} to the two
selected losses and the proposed positive alone. More generally,
with three conditional sources there is a gluing test uniform
in the number of owners. Choose at most nine minimal zeros with
consistent background union $b$. After dropping their background
postings, every internal bit $(i,d)$ impossible on these rows must
already satisfy $\pi_i(b,d)=0$ whenever $(b,d)$ is proper.
This rule is also sufficient. Indeed for three internal inputs a
repeated color in a negative occupies both passive slots, so that
row by itself is exactly one singleton-profile demand. All other
negative rows use each color at most once. The marked, rainbow and
gap certificate bounds in the proof of
\cref{core:cor:ternary-uniform-interpolation} therefore apply to
the partial internal sample, with each demanded profile costing
its one original negative. They give at most nine negatives for
an impossible internal positive. Their backgrounds are contained
in that of the positive by \cref{core:lem:unmarked-background}.
Taking their union gives the stated rule. The ternary tests and the
number of selections are polynomial even with $K$ supplied as input.

\begin{corollary}[Joint guards with components of size at most two]
\label{core:cor:binary-source-components}
For a candidate family of complete source predicates, draw an edge
$j\to i$ if changing the passive status of $j$ can change the
predicate at $i$. If every strongly connected component has one or
two inputs, joint core realizability has an intrinsic test using
only unary and binary core operations.
\end{corollary}
\begin{proof}
Check antitonicity and agreement across observers. An identically zero
source must have no outgoing dependencies: any variable occurring in
a word supplies its winner on its all-neutral source probe. Delete
such dummy inputs and reject an empty support. For prime-guard input
also check canonicity. Order the remaining components topologically.
The candidate at a prefix is reconstructed from its source predicates,
which ignore later statuses. For a core realizer this prefix must be
its neutral residual deleting the later inputs, hence is core. At each
additional component $W$ the new prefix lies between the previous
prefix $B$ and $B+U_W$, directly by source reconstruction.

Compare the first prefix with the unary or binary list of
\cref{core:prop:binary}. For each subsequent component use
\cref{core:lem:block-source-patterns} to filter the same list by
its source bits and compare the required predicates. The activation
product gives the principal boxes for the next step.
For a projection the principal tests ordinary subcoloring at its
own coordinate. Every binary operation in the list using both
inputs is its greatest profile completion, since every nonsingleton
output is input union. Its principal is therefore the explicit
profile test of \cref{core:thm:profile-source-interval}.
These finite comparisons give necessity and sufficiency at each
prefix. The enumeration on a component uses at most
$2^{K+1}+K+1$ operations, independently of the number of inputs in
the preceding components.
\end{proof}

For instance, take distinct $a,b$, let $B=x\star_a y$ and
$D=z\star_b w$, and put $F=B+(B\star_a D)$. The first pair always
contributes $B$. The second contributes $D$ exactly when neither
earlier input equals $\{a\}$. Thus $x,y$ affect one another's source
predicates, as do $z,w$, and the predicates in the second pair also
depend on the first. Here the representing coefficient is zero on
backgrounds coloring $x$ or $y$ by $a$, and $D$ elsewhere.

\paragraph{When conditional inputs occur in pairs.}
Delete unused inputs as in \cref{core:cor:binary-source-components}.
Write $U$ for the inputs retained at every coloring and $H$ for the
remaining live inputs. Thus each $i\in H$ has a minimal loss. In the
undirected source graph, two inputs are adjacent if a minimal loss
at either posts the other. For an edge $ij$ within $H$, consider the
condition
\begin{equation}\label{core:eq:reciprocal-postings}
 \text{for some }r_{ij}\in I,\quad
 \text{every loss at $i$ posts $j:r_{ij}$ and every loss at $j$ posts
 $i:r_{ij}$.}
\end{equation}
Such an owner, when it exists, is unique.

\begin{theorem}[Source guards on triangle-free graphs]
\label{core:thm:triangle-free-source-guards}
Let a live, complete source family be antitone and agree across observers.
Suppose no edge within $H$ lies in a triangle. For every such edge $ij$
and every pair of minimal losses at $i,j$ with at least one posting the
opposite input, suppose their postings outside $\{i,j\}$ together omit
an owner. Then the family is core exactly when
\eqref{core:eq:reciprocal-postings} holds on every edge within $H$.
This test is polynomial on the canonical prime lists under the stated
hypotheses.

On any graph satisfying \eqref{core:eq:reciprocal-postings}, a realizing
word is the union of the projections in $U$, the choices
$x_i\star_{r_{ij}}x_j$ on the edges within $H$, and one-source
completions on $\{i\}\cup U$, $i\in H$.
In particular, when every input is conditional the edge choices alone
give the entire operation.
\end{theorem}
\begin{proof}
Orient an edge by choosing a minimal loss $p$ at $j$ that posts
$i:r$. Take any minimal loss $q$ at $i$. Outside $i,j$, the two
neighborhoods are disjoint. The outside postings of $p,q$ therefore
have a consistent union $b$, proper by hypothesis. At $b$, with
$i$ neutral, the $j$-bit is positive: its restriction to neighbors
of $j$ is $p$ with $i:r$ deleted, and $p$ is minimal.
Apply \cref{core:lem:unmarked-background} to this positive and the
two losses, keeping just $i,j$. The binary witness uses $j$, and
$j$ alone violates $p$, so it also uses $i$. The classification
in \cref{core:prop:binary} says that the only binary operation
using and losing both sources is $x_i\star_s x_j$; it loses $j$
at $i:r$ only for $s=r$. Thus $q$ posts $j:r$. Now each minimal
loss at $i$ posts $j:r$. Fix one of them and reverse the argument
to obtain the same conclusion for every loss at $j$.

For the construction, let $e_i$ post each conditional neighbor of
$i$ at its edge owner. Every loss at $i$ contains $e_i$, and no
minimal loss can mention any other conditional input. The incident
edge choices supply $i$ exactly when a required posting is absent.
To supply its remaining positives, prescribe on $U$ the predicate
\[
       P_i(d)=1\quad\Longleftrightarrow\quad
       q|_U\not\subseteq d\ \text{for every minimal loss $q$ at $i$}.
\]
Under $e_i$ this gives precisely the desired bit wherever the
coloring is proper. Its greatest single-source completion exists
by \cref{core:thm:source-completion}; it retains $U$ and only
affects the one conditional source $i$. Omit it if $P_i=0$.
For $U=\varnothing$ every $P_i$ is zero, since $i$ has a loss.
Taking the indicated union realizes every source bit. This construction
uses neither the triangle nor the palette hypothesis. If there are
only conditional inputs, an isolated vertex would have no possible
nonempty minimal loss, so the edge sum is nonempty.

On the facet omitting $k$, every loss at $i$ posts $j:r$ exactly
when the singleton exclusion of $j:r$ suffices for a positive
$i$-bit. For $k=r$ use the empty exclusion. Sufficiency means that
a listed prime is contained in this exclusion. The edges can also
be read from the lists: each literal of a prime is necessary there,
so erasing that posting in a counterexample to the shortened guard
changes the bit. Conversely a posting that changes a bit violates
a prime satisfied before the posting. This gives the polynomial test.
\end{proof}

The palette hypothesis always holds if one owner occurs in no minimal
loss, or if $\deg(i)+\deg(j)\le K+1$ on the conditional edges: at
most $\deg(i)+\deg(j)-2$ outside inputs can be posted. In particular
the theorem covers every path and every cycle of length at least four,
at every $K\ge3$.

The degree bound is sharp even on trees. Partition $I=A\sqcup B$
with both parts nonempty, choose $a\in A$, and take a retained input
$u_r$ for each owner. Let $x$ have just the minimal loss posting
$y:a$ and $u_r:r$ for $r\in A$. Let $y$ have just the minimal loss
posting $u_r:r$ for $r\in B$. Both losses are proper.
Their outside union uses every owner, so the two-source criterion
of \cref{core:cor:two-conditional-guard-complexity} gives a core
operation. Yet $y$ affects $x$ and no loss at $y$ posts $x$.
The conditional edge has degree sum $K+2$. Retaining these same two
minimal losses after adding a fresh owner makes the family
nonrealizable, since their outside union now omits that owner.

\begin{corollary}[Source guards on a star]\label{core:thm:star-source-guards}
Let every satellite $j$ be live and depend only on the hub $x$.
Write $R_j$ for the owners at which posting $x$ suppresses $j$.
An antitone observer-agreeing family is core precisely when every
hub loss posts $j:r$ for every $j$ and every $r\in R_j$.
If every $R_j$ is nonempty and the hub has a loss, each
$R_j=\{r_j\}$ and the operation is exactly $\sum_j(x\star_{r_j}y_j)$.
\end{corollary}
\begin{proof}
For a live hub, apply \cref{core:thm:triangle-free-source-guards}:
a conditional satellite has only single-post losses, so the outside
postings in an opposite-head pair are contained in the hub loss.
A retained satellite contributes no restriction. For an unused hub
the condition holds exactly when every satellite is retained, giving
their union; any outgoing dependence from an unused hub is impossible.
\end{proof}

\paragraph{Gluing the observer facets.}
For $k\in I$, delete all $k$-postings from a coloring $q$ to obtain
$q^{\widehat k}$. The $k$-facet family extends to all proper colorings by
$\pi_i^{(k)}(q)=\pi_i(q^{\widehat k})$. Write $U_k$ for the inputs
retained everywhere on this facet. Consider the following condition
on the minimal losses $z$ at every head:
\begin{equation}\label{core:eq:facet-retained-posting}
             z_k\ne\varnothing\quad\Longrightarrow\quad
                         z_k\cap U_k\ne\varnothing .
\end{equation}

\begin{proposition}[Realization from observer facets]
\label{core:prop:facet-source-gluing}
If a family is core, each of its facet extensions $\pi^{(k)}$ is core.
Conversely, if every extension is core and
\eqref{core:eq:facet-retained-posting} holds for every $k$, the original
family is core. These implications hold on any source graph.
\end{proposition}
\begin{proof}
Read source survival with the unowned mark of
\cref{core:thm:block-source-extension}. In a realizing word replace all
$k$-owned gates by union. On a $k$-free coloring no input is the
$k$-singleton in this test, hence no subterm
takes that singleton and the replacement preserves every surviving mark.
In the resulting word, restoring a $k$-posting cannot change an
absorbing test at any gate: the remaining owners' singleton tests
depend only on their own postings. It therefore realizes $\pi^{(k)}$.

For the converse take a realizer $w_k$ of $\pi^{(k)}$ and attach
each projection in $U_k$ by $\star_k$. On the $k$-free facet no
attachment absorbs, and the extra projections are retained there.
At a minimal loss using $k$, a projection in $U_k$ is the $k$-singleton
by \eqref{core:eq:facet-retained-posting}. The entire fold then equals
that singleton and drops the marked source. Every other minimal loss
already lies on the $k$-free facet. Thus this word obeys all losses
and supplies every positive on its facet. Their union over $k$
supplies all positives since every proper coloring omits an owner.
\end{proof}

\begin{lemma}[Opposite losses over a monochromatic coalition]
\label{core:lem:monochromatic-opposite-losses}
On any source graph, take a loss $z$ at $h$, a nonempty color class
$A=z_k$, and $H=\{h\}\cup A$. At each $j\in A$ choose a loss $q^j$
whose postings within $H$ are empty or consist just of $h$ at a color
different from $k$. Let $b$ be the union outside $H$ of the postings
of $z$ and all $q^j$. If $b$ is a consistent proper coloring, every
head in $H$ must already be lost at $b$, with the other inputs in $H$
neutral.
\end{lemma}
\begin{proof}
Any contrary positive gives, by \cref{core:lem:unmarked-background},
a core word $u$ on $H$ satisfying all the indicated internal losses.
Give $h$ the menu $I$ and all of $A$ the $k$-singleton. The loss at
$h$ excludes every other winner since only $h$ could supply one.
Hence $\Phi_k(u)(A)=1$. Starting at the root, keep one $k$-winning
child at each $k$-gate and both children at other gates or union.
The singleton profile laws leave a nonempty word with all leaves
in $A$.

For each $j\in A$, probe its internal loss for observer $k$, giving
$j$ menu $I$ and the holes $I\setminus\{k\}$. This reads the same
coefficient as the canonical probe. No input is the $k$-singleton;
each deleted $k$-gate acts as union. Monotonicity therefore preserves
the omission at every one of these probes. But any $j$ occurring in
the pruned word supplies $k$ at its own probe: every input in $A$
is ambiguous there, so every gate takes union. This contradiction
proves the assertion.
\end{proof}

\begin{theorem}[Recognition on source graphs of girth at least five]
\label{core:thm:girth-five-source-guards}
Let the live, complete source family be antitone and agree across
observers. Suppose its underlying undirected source graph has no
triangles or four-cycles. It is core exactly when, for every $k$,
\begin{enumerate}
\item on the graph of the $k$-free facet, every edge between two
nonretained inputs satisfies \eqref{core:eq:reciprocal-postings}
for the minimal losses on that facet;
\item every minimal loss satisfies \eqref{core:eq:facet-retained-posting}.
\end{enumerate}
These conditions can be tested in polynomial time on complete canonical
prime lists, uniformly in $K\ge3$. In particular they decide realization
on every forest, with arbitrary directions on its edges.
When they hold, one can construct a realizing word with union named
and with $O(L)$ leaves, where $L$ counts one per head--observer list,
one per listed guard and one per literal occurrence. Its completion
depth is at most $K$, and this uniform depth bound is sharp on stars.
\end{theorem}
\begin{proof}
The minimal losses of $\pi^{(k)}$ are precisely the original minimal
losses omitting $k$. Its source graph is a subgraph of the original,
and $k$ is absent from all its losses. By
\cref{core:thm:triangle-free-source-guards}, the first condition is
equivalent to realizability of this extension. Thus
\cref{core:prop:facet-source-gluing} gives sufficiency and necessity
of the first condition.

For necessity of the second, take a minimal loss $z$ at $h$ with
$A=z_k\ne\varnothing$ and $A\cap U_k=\varnothing$. Choose a $k$-free
minimal loss $q^j$ at each $j\in A$. In a triangle-free graph none
of these posts another member of $A$ and none of its outside postings
is at a neighbor of $h$. Thus whenever their postings outside
$H=\{h\}\cup A$ are consistent, they unite with those of $z$ to
give a $k$-free coloring $b$. On the neighbors of $h$ it agrees
with $z^{\widehat k}$, where $h$ is live by minimality. This
contradicts \cref{core:lem:monochromatic-opposite-losses}.
In the absence of four-cycles the outside neighborhoods of the
distinct $j$ are disjoint, so consistency is automatic.

For the algorithm, $U_k$ consists of the heads with the empty prime
guard on the $k$-facet. Read that facet's edges from its prime literals
and apply the singleton-exclusion test of
\cref{core:thm:triangle-free-source-guards} on this list alone: a
posting is required in every $k$-free loss precisely when excluding
it suffices on this facet. No other observer's presentation of the
extension is needed.

To test the second condition, enumerate a head $h$, distinct owners
$k,\ell$, and a prime $G$ at $h$ on the $\ell$-facet. On that facet
forbid the postings $j:k$ for $j\in U_k$, and every posting in $G$
whose color differs from $k$. Test whether the allowed postings
contain a loss, using the product-domain guard scan in the proof of
\cref{core:thm:girth-five-cofactor}. Any minimal loss so obtained
must use $k$, since deleting all $k$-postings satisfies $G$.
Conversely a minimal loss violating \eqref{core:eq:facet-retained-posting}
omits some $\ell$, and deleting its nonempty $k$-block satisfies
some $\ell$-facet prime $G$ by minimality. It is found by this test.
The scans and the number of choices are polynomial in the input.

For synthesis on the $k$-facet, the losses at a conditional head $i$
all consist of its mandatory conditional-neighbor postings $e_i$
together with postings on $U_k$. The minimal hitting sets are
therefore the singleton exclusions of members of $e_i$ and the
minimal hitting sets of the remaining parts on $U_k$. If some
remaining part is empty, there are no guards of the latter kind.
Include the retained projections and the owned choice for each
conditional edge, once. These choices supply a head whenever one
of its mandatory postings is absent. For each remaining prime $G$
at $i$, start at $x_i$. For every color $r$ occurring in $G$, take
the $\star_r$-fold of the current word and the projections $x_j$ for
all $j:r$ excluded by $G$. This spine supplies the $i$-mark exactly
when $G$ is satisfied; its other leaves are retained. Thus the sum
of these terms realizes the facet extension, including heads whose
only contributions come from edges.

Each spine has $1+|G|$ leaves. Charge the two leaves of each
edge choice to a singleton prime at one of its conditional endpoints;
distinct edges have distinct charges. Each retained projection
occurs once in the facet sum and once more in the final $\star_k$-fold
of \cref{core:prop:facet-source-gluing}. This gives the stated
linear count after summing over $k$. A fold at one owner is one
greatest-completion gate: its owner wins on any nonempty coalition,
and every other owner requires unanimity. Each spine uses at most
$K-1$ such layers since $k$ is absent. An edge choice uses one
and gluing adds at most one. Free unions preserve the maximum depth.
The single-head family of \cref{core:thm:one-source-depth}, with
one distinct passive neighbor posted at each of the $K$ colors,
has a star graph and requires all $K$ layers.
\end{proof}

\paragraph{Fixing the retained inputs.}
Cycles through retained inputs can be handled without any restriction
on how many of their neighbors they control. Let $U$ be the globally
retained inputs and $W=V\setminus U$. Fix a coloring $e$ on $U$
omitting $k$. Call an original minimal loss \emph{active at $e$} if
its postings in $U$ are contained in $e$, and write $t=z|_W$ for its
internal postings. On $k$-free internal colorings put
$P_i(d)=\pi_i(e,d)$ for each $i\in W$. Write $T$ for the heads on
which $P_i$ is identically one.

\begin{theorem}[Recognition from the graph on conditional inputs]
\label{core:thm:conditional-girth-five}
Suppose the graph induced by $W$ has no triangle or four-cycle;
the whole graph may contain both. A live, antitone, observer-agreeing
family is core if and only if the following tests hold for every $e,k$
for which some internal $P_i$ is positive:
\begin{enumerate}
\item An identically zero $P_i$ has no outgoing dependence. After
deleting these inputs, every edge between two conditional heads of
the restricted family satisfies the matching rule on the $k$-free facet.
\item For each active minimal loss at $i$ with internal postings $t$,
if $P_i(t^{\widehat k})=1$, then $t_k\cap T\ne\varnothing$.
\end{enumerate}
Failure has a polynomial-size certificate from the complete prime
lists. Recognition is coNP-complete at every fixed $K\ge3$.
\end{theorem}
\begin{proof}
The facet family $P$ extends by erasing $k$-postings. Once zero
heads are removed, \cref{core:thm:triangle-free-source-guards} gives
a core realizer for this extension exactly under the first test,
since its graph lies in the graph on $W$ and its losses omit $k$.
If the original family is core, fix and erase $U$ at $e$ by
\cref{core:thm:completion-erasure}. The residual is nonzero since it
supplies a positive; on $k$-free colorings it has the stated bits.
Replacing its $k$-gates by union gives the required facet extension.
It also follows by \cref{core:lem:unmarked-background}, applied to
each positive and then taking their union, that one word on $W$
supplies all these positives while omitting all active internal losses.

Take an active loss with $P_i(t^{\widehat k})=1$.
The set $A=t_k$ is nonempty, since the active loss itself forces
$P_i(t)=0$ if $t$ omits $k$. If $A$ misses $T$, for each $j\in A$
choose a $k$-free zero of $P_j$. It contains the internal part of
an active $k$-free original minimal loss: restrict a minimal
original loss inside the fixing by $e$ and that zero.
In the graph on $W$, these internal losses post only neighbors of
their respective heads. Their postings outside $\{i\}\cup A$ are
pairwise disjoint and avoid all neighbors of $i$ by the girth
hypothesis. The $i$-bit is therefore still positive when they are
added to $t^{\widehat k}$, with $A$ neutral. This contradicts
\cref{core:lem:monochromatic-opposite-losses} for the internal word.

Conversely, start with the facet realizer from the first test and
attach every projection in $T$ by $\star_k$. It still realizes
every $k$-free positive and fits every active internal loss.
Indeed, if $P_i(t^{\widehat k})=0$, the facet realizer already omits
the mark there and $i\notin T$. Otherwise a member of $T$ posted
at $k$ absorbs the fold. The background lemma now lifts any requested
positive with that fixing. The globally retained projections and a
union over the internal positives realize the family.

To compute the complete restricted primes, drop each $k$-facet
guard violated by $e$, remove its $U$-literals from each surviving
guard and take the inclusion-minimal results. For completeness,
adjoin to a sufficient restricted guard all exclusions on $U$ not
violated by $e$. The resulting guard is sound before restriction:
on $U$ it allows only subcolorings of $e$. It therefore contains an
original prime which survives at $e$. This proves that the first
test and $T$ are computable in polynomial time for specified $e,k$.
For failure of the second, exhibit one original minimal loss and
evaluate its restricted bit; minimality is checked by single-post
deletions. If $W$ is empty the retained projections suffice.
The lower bound is \cref{core:cor:two-by-n-source-hardness}, where
the conditional vertices induce a three-vertex path.
\end{proof}

With four-cycles, the losses at the members of $A=z_k$ may impose
conflicting postings on a shared input. For a minimal loss $z$ at
$h$ and each nonempty $A=z_k$, impose the following rule:
\begin{equation}\label{core:eq:compatible-facet-losses}
 \begin{gathered}
 \text{there is no choice of one $k$-free minimal loss at each
 $j\in A$}\\
 \text{whose postings outside $\{h\}\cup A$ agree wherever they overlap.}
 \end{gathered}
\end{equation}

\begin{theorem}[Triangle-free dependence with shared retained inputs]
\label{core:thm:triangle-free-shared-inputs}
For a live, antitone, observer-agreeing family with triangle-free
source graph, facetwise matching and
\eqref{core:eq:compatible-facet-losses} are necessary for realization.
Together they imply full colored-cofactor independence. They characterize
realization if either the dependencies are directed acyclically or every
undirected four-cycle contains an input retained everywhere.
On both classes, failure has a polynomial-size certificate from the
complete prime lists.
\end{theorem}
\begin{proof}
The necessity argument in \cref{core:thm:girth-five-source-guards}
uses absence of four-cycles only to guarantee that the outside
postings agree. For the cofactor law, fix a minimal loss $z$ at $h$
and test any winner $k$. If $z_k=\varnothing$, the $k$-facet extension
has the same loss and is core by matching. Its cofactor is independent
of $h$, and its $k$-membership equals that of the original operation
on all valuations since that observer ignores $k$-postings. Otherwise
all $j\in z_k$ carry $k$ at every valuation extending the fixing.
If none supplied it, their zero bits on the $k$-facet would give a
selection of minimal losses compatible with that valuation, violating
\eqref{core:eq:compatible-facet-losses}. Thus $k$ is always present.
Every fixing at a loss extends a minimal one. The acyclic implication
is \cref{core:thm:acyclic-source-guards}; a live input with no
predecessors ensures nonemptiness in that case.

If every square meets the globally retained set $U$, apply
\cref{core:thm:conditional-girth-five}. Its first test holds by
fixing $e$ in the core $k$-facet extension and erasing $U$.
For its second test, suppose an active $k$-block $A$ misses $T$.
Choose active $k$-free losses for the heads in $A$ as in that theorem.
Their postings on $U$ are compatible because all lie in $e$.
Away from $U$ and the original head they have disjoint supports,
since no square lies outside $U$. This contradicts
\eqref{core:eq:compatible-facet-losses}.

Facet matching is checked directly from the prime lists. For the
other rule a certificate lists $k$, one minimal loss $z$ and the
compatible losses at its $k$-posted inputs. Each listed coloring
has at most one posting per input; evaluating its prime predicate
and erasing each posting in turn verifies minimality.
\end{proof}

\begin{proposition}[One square separates the facet tests from realizability]
\label{core:prop:square-facet-obstruction}
Take three owners $a,b,k$ and retain $p,q$. Give $u,v,h$ the losses
\[
  (w:a,p:b)\longrightarrow u,\qquad
  (w:b,q:a)\longrightarrow v,\qquad
  (u:k,v:k)\longrightarrow h .
\]
Retaining $w$ gives a core operation. Giving $w$ the single minimal
loss $(u:a,v:b)\longrightarrow w$ gives a noncore operation.
Both families satisfy the two conditions of
\cref{core:thm:triangle-free-shared-inputs}, and their undirected graph
is the same four-cycle with two pendant vertices.
\end{proposition}
\begin{proof}
In both families, the $a$- and $b$-free facets have only $h$
conditional. On the $k$-free facet with $w$ retained only $u,v$
are conditional, and neither affects the other. With the additional
loss, the $k$-free conditional edges $uw,vw$ match at $a,b$,
respectively. In \eqref{core:eq:compatible-facet-losses}, the
$k$-postings at the $h$-loss are protected because losses at $u,v$
require different colors at $w$. Each color class of every other
listed loss contains an input retained on that color's facet.
The first family is directed acyclic, so the corollary realizes it.

For the second family ask for the positive at $h$ on the coloring
$P=(u:a,v:b,p:b,q:a)$, with $w$ neutral. Give a marked head the
ambiguous menu $I$ as above. A word supplying the mark cannot have
a winning $a$-test on $P_a$ or a winning $b$-test on $P_b$.
It therefore cannot be singleton on the $w$-loss. On the $u$-loss
its only possible singleton is $a$, since $p\in P_b$.
That would give $\Phi_a(\{w\})=1$. On the $w$-loss the word omits
the mark carried by $w$, so singleton-profile absorption
(\cref{core:lem:singleton-profile-absorption}) forces the value
$\{a\}$ there. This gives $\Phi_a(\{u\})=1$, impossible since
$u\in P_a$. The $v$-loss is symmetric.

A singleton on the $h$-loss would give $\Phi_k(\{u,v\})=1$.
Prune to one winning child at each $k$-gate as in
\cref{core:lem:monochromatic-opposite-losses}. The resulting word uses
only $u,v$ and still omits the mark on both their $k$-free loss probes.
Both of its inputs are ambiguous on each of those probes, so any leaf
supplies the forbidden mark on its own row. This rules out the last
singleton. The root must therefore take union on all four losses.
Both children omit those marks and one still supplies the positive.
Descent would reach the projection $h$, which fails its own loss.
\end{proof}

\begin{corollary}[Two common neighbors suffice for hardness]
\label{core:cor:two-by-n-source-hardness}
For each fixed $K\ge3$, recognition from complete canonical prime
lists is coNP-complete on acyclic orientations of complete bipartite
graphs $K_{2,n}$ of depth two. Deleting either vertex of the two-element
side leaves a tree. In particular both classes of
\cref{core:thm:triangle-free-shared-inputs} have coNP-complete recognition.
\end{corollary}
\begin{proof}
Use the hypergraph two-colorability instances of
\cref{core:cor:two-conditional-guard-complexity}, deleting isolated
vertices and taking the edges to be a nonempty antichain of nonempty
sets. Choose distinct owners $a,b,k$. Retain the vertex inputs and
one flag $f_r$ for each $r\ne k$. Let $u$ lose when all flags have
their indicated colors and the vertices posted at $a$ meet every edge.
Use the same rule at $v$ with $b$ in place of $a$. Let $h$ lose
exactly when $u,v$ are both posted at $k$. The two-element side
is $\{u,v\}$. Each vertex input affects both predicates: for an
edge containing that vertex one can post the vertex and everything
outside that edge, meeting every edge by the antichain condition,
then erase the vertex to miss the chosen edge.

On the $k$-facet the primes for $u$ exclude one flag posting or
all the $a$-postings of one edge; all allowed $a$-postings can be
made simultaneously. This proves sufficiency and exhaustiveness,
with minimality by the antichain assumption. The $v$-list is analogous.
Other facets of $u,v$ have the empty guard. At $h$ the only nontrivial
primes exclude $u:k$ or $v:k$. These complete lists have polynomial size.

Two disjoint transversals extend to a proper two-coloring of the
hypergraph, and the two color classes of any such coloring are
transversals. If they exist, post them at $a,b$, set every flag
and fix $u,v$ to $k$. Neither intermediate input can now supply $k$,
and changing $h$ to $\{k\}$ introduces it, contrary to the cofactor
law. If there are no disjoint transversals, at least one of $u,v$
always supplies $k$ under this fixing. The $h$-bit is zero at every
other observer there. Under a fixing suppressing $u$ or $v$ the flags
supply all winners except $k$, and the $k$-bit of $h$ is unconditional.
These are all cofactor tests, giving realizability by
\cref{core:thm:acyclic-source-guards}. Membership is the acyclic upper
bound of \cref{core:thm:complete-acyclic-guard-complexity}.
\end{proof}

If both the prime guards and the minimal losses are given, recognition
on these acyclic $K_{2,n}$ graphs takes polynomial time. More generally
this holds for an acyclic dependency graph supplied with a tree
decomposition of fixed width $w$. To find a cofactor violation,
enumerate its captured head $h$ and one minimal loss there, the observed
winner $k$, the two menu classes at $h$, a positive supplier and one
of its $k$-primes on the second valuation. As in
\cref{core:thm:girth-five-cofactor}, this imposes only individual menu
restrictions. On the first valuation each input $i$ must either omit
$k$ or satisfy one of its listed $k$-free losses. Introduce a variable
for this choice at $i$. Each choice places individual restrictions
on the menu at $i$ or at its predecessors. This is a binary constraint
problem, with at most $K+2$ menu classes and one more choice than
the number of listed losses per input. Replace each input in each
bag by its menu and choice variables. Every constraint then lies in
a bag of size at most $2w+2$. Tabulate the assignments in each bag
satisfying its constraints and matching a feasible assignment in each
child on their intersection. Connectedness of the bags for each variable
makes this test exact, and its cost is polynomial at fixed $w$.
For $K_{2,n}$ one bag per vertex on the large side, together with
the two common vertices, gives width two.

\paragraph{One positive coloring through the source components.}
Consider a sample of proper negative rows and order the strongly
connected components of its posting-to-head graph topologically. Write
a row at $i\in C$ as $(i,b,t)$, where $t$ gives the postings inside
$C$ and $b$ the earlier postings. Fix a proper coloring $p$ of the
inputs. The principal $\Theta_{0,q}$ for the absent operation is
constantly true.

\begin{theorem}[A component pass for partial source queries]
\label{core:thm:component-mark-sweep}
Process the components in order. On $C$ keep the internal loss
$(i,t)$ from each row $(i,b,t)$ satisfying
\[
       \Theta_{D_U,p|_U}(b|_U)=1
       \qquad\text{at every earlier component }U.
\]
Let $D_C$ be the union of all core operations on $C$ omitting the
kept internal losses, or zero if none does. A word omitting the
original sample supplies the mark at $i\in C$ under
$p|_{V\setminus\{i\}}$ exactly when
$\pi_i^{D_C}(p|_{C\setminus\{i\}})=1$.

For a listed sample the cost is a function of $K$ and the largest
component size $w$ times a polynomial in the input length. At this
$p$, all the answers are preserved by a common subfamily of at most
$\sum_C |C|(K+1)^{|C|-1}$ original rows.
\end{theorem}
\begin{proof}
At each prefix take the union $B$ of all words on that prefix
omitting its rows, allowing zero. The earlier source bits of any
bounded word on a larger prefix lie below those of $B$, even if
later inputs are posted: a witness for such a bit localizes to
the source's ancestors by \cref{core:thm:component-profile-factorization}.
Conversely $B$, when nonzero, omits all later-headed marks.

For an extension by $C$ with $B\ne0$, take the upper bound whose
earlier source predicates are those of $B$ and whose new source
predicates lose just on extensions of the given new rows.
Source reconstruction puts it between $B$ and $B+U_C$.
By the preceding observation its core minorants are exactly
the bounded words. By \cref{core:lem:block-source-patterns}, the
coefficient of their union at the earlier part of $p$ is the
greatest local word omitting those internal losses for which
$b\preceq_B p|_{\mathrm{prefix}}$. The activation product of
\cref{core:thm:block-source-extension} identifies this comparison
with the displayed tests, and multiplies it by the factor
$\Theta_{D_C,p|_C}$ for the next step. This uses each original
proper row: its internal postings need not give a proper coloring
when combined with the positive background.

If $B=0$, no word on the extended prefix satisfying the rows can
use an earlier input. Any occurrence supplies its own mark at the
empty coloring, which would localize to the earlier prefix. Hence
all earlier postings can be erased and every earlier factor is
constantly true. This also starts the induction at the empty prefix.

For fixed $K,w$ enumerate the local term tables and their principals
once, then test the displayed conditions directly. To preserve all
$D_C$ keep one original row for each different internal loss retained
in the pass. The new pass has identical local losses and coefficients
by induction. There are at most $(K+1)^{|C|-1}$ patterns per head.
\end{proof}

\begin{corollary}[One internal posting owner per component]
\label{core:cor:component-owner-marks}
Suppose all postings from a component $C$ to its own heads have
the same owner $r_C$. A specified marked query takes polynomial
time on listed negatives, for arbitrary positive colorings and
component sizes. At the uncolored query the same holds on complete
canonical prime lists, without listing minimal zeros.
\end{corollary}
\begin{proof}
By \cref{core:thm:conditional-component-capture}, each retained
loss at $i$ with internal tail $A$ gives the Horn rule $A\to i$.
For their closure $c$, the coefficient is a one-owner trace on
$C\setminus c(\varnothing)$, or zero if this is empty.
\Cref{core:thm:one-owner-activation-order} evaluates every nonzero
coefficient's comparison by Horn closure on this support. Its
source $i$ supplies the mark at $p$ exactly when
$i\notin c((p_{r_C}\cap C)\setminus\{i\})$. Thus the component pass on a listed
sample needs polynomially many closure computations.

At the uncolored query only the live support of each coefficient is
needed: an earlier comparison allows precisely the colorings leaving
that support unposted, by \cref{core:thm:activation-faithfulness}.
To compute $c(\varnothing)$ from primes, maintain the derived set $X$.
For each underived head $i$ allow internal postings only at $r_C$
on $X$, forbid postings on earlier live supports and on later inputs,
and leave the remaining earlier inputs unrestricted. On each facet
these exclusions form a guard. Some retained rule derives $i$
exactly when one such guard contains no listed prime, since its box
then contains a zero. Forward chaining therefore computes every
live support in polynomial time.
\end{proof}

\begin{theorem}[Prime-guard complexity from source components]
\label{core:thm:prime-component-space}
Fix $K\ge3$ and let $N$ be the length of the complete canonical
prime lists of an observer-agreeing antitone family. At bounded
component size $w$, feasibility of a specified positive query is in
coNP. It is polynomial for $w=1$ and coNP-complete already for
disjoint two-cycles feeding one additional source, even with at
most three exclusions per prime.

If $c$ inputs belong to nontrivial components meeting the support of
the positive coloring, the specified query takes time $f(K,w,c)N^{O(1)}$.
In particular the uncolored query is polynomial at bounded $w$.
For full realizability the complexity at bounded $w$ is coNP-complete
already for acyclic dependence. At arbitrary $w$ full recognition uses
deterministic space $N^{O(1)}2^{O_K(w)}$, hence polynomial space for
$w=O(\log N)$.
\end{theorem}
\begin{proof}
A posting belongs to a minimal zero on a facet exactly when it
occurs in a prime, by minimal hitting-set duality. Thus the lists
give the graph. For a failed query guess the subfamily in
\cref{core:thm:component-mark-sweep}, check that each selected row
is zero on a missing owner's facet, and run the explicit pass.
For full recognition also guess the missing positive: a common word
exists precisely when the individual positives can all be supplied,
by union, provided there is at least one positive. A family with
no positive bit has no realizer. The acyclic bounds for the two
problems are \cref{core:thm:acyclic-marked-source-test,core:thm:complete-acyclic-guard-complexity}.

For fixed-query hardness reduce 3CNF unsatisfiability. Remove
redundant clauses containing another, taking nonempty clauses.
Give each variable two inputs for its opposite literals, with a
loss at each when the other is posted at $a$. Their only bounded
binary operation is $\star_a$ by \cref{core:prop:binary}. At the
positive post the first input of every pair at $a$. Its local
principal admits exactly the colorings that do not put both
literals at the same foreign color: this is $\theta_a$ of
\cref{core:thm:profile-source-interval}, since the $a$-profile is
disjunction. At the final head $h$ prescribe loss when the literals
posted at $b\ne a$ meet every clause. On facets seeing $b$, the
canonical primes exclude all the $b$-postings of one clause;
the clause antichain makes these exactly the minimal such guards.
The other facet has the empty guard. A loss passes the pair
principals precisely when a consistent choice of literals meets
every clause: a satisfying assignment uses only $b$-postings, and
an allowed coloring cannot post opposite literals at $b$. Thus the
$h$-mark is feasible exactly for an unsatisfiable formula.

For the time bound, an earlier singleton block permits every posting
if $D_U=0$, and otherwise only subcolorings of $p|_U$. An uncolored
block permits exactly the colorings with no post on the support of
$D_U$: each used source supplies its all-neutral mark, so the
predicate $Q_j$ of \cref{core:thm:activation-faithfulness} excludes
any post there; unused inputs have no effect. Enumerate the colorings
on the remaining earlier blocks, involving at most $c$ inputs, and
discard choices failing their local principals. Also enumerate the
internal coloring $t$ at the head being tested. To test whether the
local loss $t$ or a subcoloring is forced, allow every subcoloring
of these choices, impose the indicated individual restrictions on
earlier singleton and uncolored blocks, and leave later inputs
unposted. On each observer facet these restrictions enable one
guard. It contains a listed prime exactly when its entire box is
positive. A zero in a tested box gives a retained internal loss
below $t$, by antitonicity of the principals. Conversely every
retained loss is tested at an observer absent from its coloring.
Discard a candidate local word if any of its positive internal
patterns admits such a loss. This computes $D_C$ within the claimed
bound.

For the general space bound, stream all positive probes through the
pass. Each local table has at most $(2^K-1)^w$ entries. Its core
membership is a subpower test on the explicit table, using space
polynomial in its length by \cref{core:thm:subpower-complexity}.
Compute $p'\preceq_D q$ on a local table $D$ by enumerating core
tables above $D$ and below $D+X$ on one extra variable: none may
separate the two bits in the forbidden direction. Zero allows every
comparison. Testing a trial local word requires only a scan of
original zero patterns and the already stored local tables.
Storing the union of the successful tables at each prefix and this
workspace costs $N^{O(1)}2^{O_K(w)}$ space. If agreement or canonicity
is to be checked, scanning all patterns and candidate guards also
uses polynomial space.
\end{proof}

\begin{corollary}[Distinct posting colors within each component]
\label{core:cor:component-injective-marks}
Suppose the postings from the head's own component have distinct
colors in every negative row. A specified marked query on explicitly
listed negatives takes polynomial time, with no bound on the sizes
or number of components or on the repetitions in earlier components.
At fixed $K\ge3$, feasibility from canonical prime lists is coNP-complete
under this restriction on the minimal zeros, as is full realization.
The specified uncolored query takes polynomial time even on the
prime presentation.
\end{corollary}
\begin{proof}
In the component pass there is no need to construct the table $D$
of a component $C$. Let $S$ be its retained internal losses. For
proper $b,p$ on $C$, the comparison $b\preceq_D p$ fails exactly
when a word omitting $S$ supplies a fresh source $Y$ at $p$ while
omitting it at $b$. Indeed an interval separator above $D$ is such a
word. Conversely any witness $g$ lies below $D+Y$: each old source
bit it supplies localizes to $C$, since no old loss posts $Y$, so
the old contribution is bounded by $D$. Thus $D+g$ is an interval
separator with the same $Y$-bits. For $D=0$ no witness can use an
old input, by applying localization to its all-neutral mark; the
projection $Y$ fails the new loss. This gives exactly the
constantly true comparison at zero.

Each comparison is therefore a marked test on $S$ with just one
extra loss, at $Y$. Only that row can repeat a posting color, so
\cref{core:thm:repeated-blocker-budgets} gives polynomial time.
The final internal mark test has no repeated class. With listed
negatives the pass makes polynomially many of these calls.

At fixed $K$ only $O_K(n^K)$ internal rows can occur across all
components: each has at most $K-1$ postings. One original row
per retained internal pattern again certifies failure, giving the
prime-input coNP bound for both problems. The marked lower bound
in \cref{core:thm:prime-component-space} uses only singleton internal
postings, and the full lower bound is already acyclic.

At the uncolored query, determine the support of every $D$ by the
internal uncolored mark tests. Earlier comparisons then merely
exclude postings on these supports. For each possible injective
internal pattern $t$, allow its subcolorings and these earlier
postings, leaving later inputs unposted. As in the prime test above,
one guard per facet decides whether the box contains an original
zero. Keeping $t$ whenever it contains a retained internal loss
gives the same bounded words. Enumerating the $O_K(n^K)$ patterns
and using the explicit local test completes the pass.
\end{proof}

\begin{corollary}[Localizing monochromatic activation to one component]
\label{core:cor:unanimity-source-component}
Let $F$ have unanimous $a$-profile and let $q$ color every input by
$a$. If $\Theta_{F,q}$ is not ordinary subcoloring, some strongly
connected component of the source graph contains at least four inputs
at which every one-input coloring is admitted. In particular, among
counterexamples of least arity at any fixed $K$, the source graph is
strongly connected and every input admits every singleton color.
\end{corollary}
\begin{proof}
At a split into a nonempty topological prefix and a new component,
the prefix $B$ also has unanimous $a$-profile. Otherwise give the
new inputs $\{a\}$ at a nonunanimous background with $B(u)=\{a\}$:
the upper bound would force the full output to be $\{a\}$.
At the all-$a$ background the block formula returns
$\{a\}\cup D_{q_J}(w)$. Unanimity on the full input makes $D_{q_J}$
nonzero, with unanimous $a$-profile. By the activation product and
induction, the principal of $F$ is the product of the principals of
these unanimous operations on their respective components.
Any nonordinary factor has at least four inputs admitting every
one-input coloring by \cref{core:cor:unanimity-released-inputs};
these colorings also pass the product. If the source graph has several
components, that factor gives a smaller counterexample. If any input
excludes a singleton color, the same cited corollary reduces arity.
\end{proof}

\paragraph{How many colors a blocking pattern needs.}
Any zero of a source predicate is visible at a valuation with at most
$K-1$ singleton colors. Choose an observer $\ell$ missing from the
pattern, give its source an ambiguous menu containing $\ell$, and give
every uncolored passive input $I\setminus\{\ell\}$. The source alone
could supply $\ell$. Consequently agreement with union on all such
valuations forces a source-reconstructed operation to be union. One
fewer color leaves the full recognition problem.

\begin{theorem}[Reduction to blocking patterns using every color but one]
\label{core:thm:mixed-guard-reduction}
Fix $K\geq3$ and $m\geq1$. From any nonempty-valued candidate $G$ on
$n\geq1$ inputs reconstructed from antitone, observer-independent
source predicates, one can form a candidate on $mK(n+K-1)$ inputs
over the same owners with the following properties.
\begin{enumerate}
\item Every departure from input union changes $I$ to
$I\setminus\{\ell\}$ for a single winner $\ell$. Each remaining
winner then occurs at least $m$ times as a singleton input. In
particular all singleton profiles are unanimous.
\item For every $O\subseteq I$, one core word can reproduce
membership of every winner in $O$ in the new candidate if and only
if one can do so in $G$. In particular it is core if and only if
$G$ is core.
\end{enumerate}
\end{theorem}
\begin{proof}
First take $m=1$. For each observer $k$, take a copy $W$ of the
original inputs. Write
$U$ for their union and $A_k$ for their $\star_k$-fold. For each
$s\ne k$ adjoin an input $X_s$. Put
$B=\sum_{s\ne k}X_s$. The test at $s$ is $X_s=\{s\}$. Define $F^k$
by adjoining the indicated entry to $B$.
\[
\begin{array}{c|ccc}
 \text{number of tests failed}&0&1&\text{at least }2\\ \hline
 \text{contribution from }W& G&G+A_k&U
\end{array}
\]
Take the union of the $F^k$ on disjoint copies of all their inputs.
Every added input contributes unconditionally. For a source in $W$ the three
entries give its predicate according to the tests on the
observer-erased background. Erasing the observed color can change
a test only when $B$ already supplies that winner. The entries
decrease as more tests pass.

If all tests pass, every color other than $k$ is present in $B$
as an input atom. If just the test at $\ell$ fails, a loss requires
an internal $k$-atom, since otherwise $A_k=U$. Then $A_k$ supplies
$k$, and $B$ supplies all other colors except possibly $\ell$.
In either case a lost winner is the only one missing and each
remaining winner occurs as an input atom. A loss in the disjoint
union must occur in a copy containing the lost winner in its inputs.
This proves the first assertion for $m=1$, including unanimity since
$K-1\geq2$.

If a core word $g$ realizes $G$, a word for $F^k$ is
\[
 B+g+\sum_{s\ne k}(A_k\star_s X_s)
 +\sum_{\substack{s<t\\s,t\ne k}}((U\star_s X_s)\star_t X_t),
\]
in any fixed order of the owners. When a test fails, its gate retains
every winner of its other operand. When it passes, that gate
contributes only its tested atom, already in $B$. This gives the
three entries. Since $g$ occurs only as a union summand, the same
construction works if $g$ agrees with $G$ only on observers in $O$.

Let $q^k$ color each added input by its prescribed color. Its
coefficient in $F^k$ is the greatest core minorant of $G$, or zero.
Indeed, any such minorant is bounded by $F^k$ on every background.
For the reverse bound, take a zero $t$ of an internal source in $G$.
If $t$ omits $k$, the joint pattern $(q^k,t)$ is proper and blocks
that source. Otherwise choose $\ell$ missing from $t$ and erase just
$X_\ell$ from $q^k$. On this subcoloring exactly
one test fails. The source is still blocked, since $G$ blocks it and
the $k$-atom in $t$ makes $A_k$ absorb. The joint pattern now
omits $\ell$. In either case \cref{core:lem:block-source-patterns}
forces the coefficient to respect the zero at $t$.

For the converse take a word agreeing with the disjoint union on
the observers in nonempty $O$. For each $k\in O$, neutrally erase
the other copies. This preserves agreement on $O$ in the $k$-copy,
since erased menus $I\setminus\{\ell\}$ cannot supply the observer
$\ell$. Erase the background at $q^k$ by the marked-word construction
in \cref{core:thm:block-source-extension}. The resulting word $h_k$
on $W$, or zero, has exactly the formal source bits of the word at
that background. On patterns omitting $k$ they are the bits of $G$.
For any zero $t$ visible to some $\ell\in O$, the preceding proof
exhibits a subcoloring of $q^k$ where the bit is zero as observed
by $k$ or $\ell$. A formal source bit is antitone even when all
colors occur, so it is still zero at $(q^k,t)$. Hence $h_k$ respects
every zero visible to $O$. Taking the union of the nonzero $h_k$
supplies every positive bit visible to $O$, and source reconstruction
gives agreement on $O$. Nonemptiness on a diagonal $\{k\}$ ensures
a nonzero summand. For $O=\varnothing$ any projection suffices.

Finally replace every input by the union of $m$ fresh inputs.
Substitution preserves realizability on each $O$ in both directions,
since identifying the $m$ copies recovers the original operation.
Such a union is the $s$-atom exactly when all its inputs are the
$s$-atom. Each of the $K-1$ colors needed for a loss now requires
$m$ singleton inputs, and the singleton profiles remain unanimous.
\end{proof}
